\section{Introdução}

O Problema do Caminho Hamiltoniano é um dos problemas mais conhecidos da ciência da computação, inclusive sendo um dos problemas NP-Completos de Karp \cite{Karp1972}. Do ponto de vista histórico, o conceito de caminho hamiltoniano remonta ao século XIX, vindo de um jogo criado por William Rowan Hamilton, e desde então tornou-se um objeto de estudo clássico na teoria dos grafos \cite{West2000}. Do ponto de vista da ciência da computação, sua relevância vem de seu valor teórico, como exemplo de problema intratável no pior caso, e também de suas relações com outros problemas de otimização e sequenciamento, como o Problema do Caixeiro Viajante, planejamento de rotas, ordenação de tarefas e análise de redes.

Seja $G = (V, E)$ um grafo não direcionado. Um caminho simples $P \in G$ é um subgrafo definido por $P = (V', E')$, onde $V' = \{v_1, v_2, \ldots,v_n\}$, $E' = \{v_1v_2, v_2v_3, \ldots,$ $v_{n-1}v_n\}$ e, $\forall v_i, v_j \in V'$, $i \neq j \Rightarrow v_i \neq v_j$. De forma mais intuitiva, um caminho simples corresponde a uma trajetória no grafo que se inicia em um vértice, percorre uma sequência de vértices adjacentes e termina em um vértice final, sem repetir vértices ao longo do percurso. Se um grafo $G$ possui um caminho simples $P$ que contém todos os seus vértices, então $P$ é chamado de Caminho Hamiltoniano.

Determinar a existência de um caminho dessa natureza em um grafo qualquer é um problema computacionalmente difícil. Mais precisamente, na sua versão de decisão, dado um grafo não direcionado $G = (V,E)$, deseja-se decidir se existe um caminho hamiltoniano em $G$. Isso define o Problema do Caminho Hamiltoniano (HAMPATH), que é conhecido por ser NP-completo.

Neste trabalho, estuda-se o Problema do Caminho Hamiltoniano de forma teórica e experimental. Teoricamente, são apresentadas as definições formais, a prova de pertencimento do problema à classe NP e sua NP-completude por meio de reduções a partir do 3-SAT para o Problema do Caminho Hamiltoniano Dirigido (DHAMPATH), que por sua vez é reduzido para HAMPATH. Também são discutidos métodos para resolução do problema, desde abordagens exatas até heurísticas e aproximações.

Experimentalmente, são comparadas duas abordagens exatas para a resolução do problema: um algoritmo baseado em programação dinâmica (Bellman-Held-Karp - BMK) e uma abordagem baseada em redução para SAT resolvida por um solucionador CDCL (\textit{Conflict-Driven Clause Learning}). O objetivo é analisar empiricamente o desempenho dessas abordagens em termos de tempo de execução e consumo de memória, discutindo suas vantagens e limitações.
